Very thorough and perceptive analysis.  I agree with your suggestion to begin with the weakest derivation of the series and make sure we have that solid before handling the details of the other papers and their already reasonable handling of the other aspects of QM.  Yes, please proceed to develop the Bell's theorem paper.

2:08 AM

Architected rigorous Bell's theorem framework for CPP theory
Architected rigorous Bell's theorem framework for CPP theory
Before writing a line, let me work through the mathematics carefully to make sure the proof is correct.


Ran 2 commands
Ran 2 commands
All mathematics verified. Now writing the paper.


Ran 3 commands, created a file, read a file
Ran 3 commands, created a file, read a file
689 lines, clean structure with four formal theorems/propositions, all proofs complete. Here is a plain account of what the paper does and why it is correct.

What the paper establishes
The error in v1 is identified explicitly and formally. Section 1 opens by naming the mistake: the shared-bit-pool argument is structurally identical to EPR local realism. Bell's theorem says this class of model can never exceed |CHSH| = 2. The v1 paper claims |CHSH| = 2√2. These are logically incompatible, and the paper says so plainly.

Two things work together to produce Bell violation. Neither one alone is sufficient:

Non-separability: The singlet DI-bit state |Ψ⁻⟩ = (|↑↓⟩ − |↓↑⟩)/√2 cannot be factored into independent states for A and B. This is Proposition 2.1, proved algebraically. Bell's theorem assumes the hidden variable factorizes — that assumption fails for a non-separable state. The Nexus maintains this non-separability as the particles separate.

Born rule: The probability rule P = |projection|² (from DI-bit density, companion C3) gives E = −cos(θ). The classical rule P ∝ |projection| gives E = 1 − 2θ/π. Appendix A works through the classical calculation explicitly and shows it saturates at |CHSH| = 2. The Born rule is what makes the difference.

Both theorems are proved from scratch. Theorem 4.2 derives |CHSH| = 2√2 from the correlation function by direct arithmetic. Theorem 5.1 proves no-signaling by showing Alice's marginal P(A=+1) = 1/2 independently of Bob's axis.

The comparison table (§6) places CPP honestly alongside Copenhagen, Many-Worlds, and Bohmian mechanics. CPP is most similar to objective collapse models but derives its collapse mechanism (bit-sea decoherence) from first principles already in the framework, not as a new postulate.

The falsifiable prediction is a lattice correction to the correlation function of order (ν/ν_P)² near the Planck frequency, plus a gravity-entanglement coupling prediction near black hole horizons. Both are beyond current experiments but are genuine CPP-specific predictions that distinguish it from standard QM.
